\section{Introduction}
\label{sec:introduction}

At a 150-year horizon, almost nobody alive at the start is alive at the end. A
projection to \ProjectionEndYear{} from a \PhaseFiveBaseYear{} base year is
therefore barely a statement about today's population at all. It is a statement
about the process that produces births, run four or five times over, and
everything else washes out. Freeze today's fertility rates and the arithmetic
gives about \ConstantFertilityTwentyOneFifty{} billion people; use the United
Nations' medium assumptions and it gives about
\UnEquivalentTwentyOneFifty{} billion. Because both calculations use the same
base data and projection engine, the difference arises entirely from the assumed
fertility path.

Conventional long-run projections make that assumption in a specific form.
Total fertility is treated as a national time series that walks down through the
demographic transition and then reverts probabilistically toward a long-run
level \citep{alkema2011,raftery2012,sevcikova2016}. The United Nations
\citep{un2024methodology} and the Institute for Health Metrics and Evaluation
\citep{vollset2020} disagree about the level but agree about the object being
modelled: one number per country per year, evolving by its own history.

That object has no composition in it, and composition is not a detail here.
Women differ a great deal in how many children they have. Across
\CVCountries{} low-fertility countries with completed cohort data the
coefficient of variation of completed family size has a median of \CVMedian{},
so the standard deviation is more than half the mean. This variation persists
across generations: adult children of larger families have somewhat larger
families themselves, with
reported parent--child correlations usually between 0.1 and 0.2
\citep{murphy1999,murphy2013}.

Together, dispersion and intergenerational persistence imply a process that a
national-rate model cannot express. The parents of the next generation are not a
random sample of this one
--- they are sampled in proportion to how many children they have. Whatever
inclined a woman toward four children rather than one is over-represented among
the people raising the next cohort, and if it carries at all, the composition of
the population moves toward it. This is selection in the sense used since
\citet{fisher1930}, and its direction on fertility is always upward, because
fertility here is simultaneously the trait and the measure of reproductive
success. A model that projects only the national mean does not estimate this
force to be small. It has no way to see it.

\subsection{What the existing arguments get right, and where they stop}
\label{sec:prior-work}

That selection acts on fertility is not a new observation. \citet{kolk2014} show
that intergenerational fertility correlations can prevent low fertility from
persisting; \citet{burger2016} argue that treating the fertility decline as
permanent is unjustifiable on evolutionary grounds; and \citet{collins2019} apply
the breeder's equation to national fertility and conclude that global population
stabilisation this century is unlikely. There is also direct evidence that
contemporary human populations are under measurable selection on
fertility-linked traits \citep{byars2010,beauchamp2016,kong2017}.

The sharpest objection to that work is also the most useful.
\citet{arenberg2022} point out that a model built from a transmission
coefficient alone contains no absolute fertility level: it can describe which
dispositions become more common without being able to say whether the population
grows, and the empirical fact that higher-fertility subgroups are themselves
drifting toward replacement has nowhere to enter. Their condition is clean --- a
high-fertility type shrinks unless the expected number of same-type daughters
each woman leaves exceeds one --- and it simply cannot be evaluated in a model
that carries no fertility schedules.

\subsection{What this paper does}
\label{sec:contribution}

\begin{enumerate}[label=\arabic*., leftmargin=1.6em]
  \item \textbf{Puts the mechanism inside a real projection.} Selection is added
  as an extra axis on a single-year, two-sex cohort-component engine for 237
  countries carrying the United Nations' own age-specific fertility, survival and
  migration schedules (\Cref{sec:engine}). Absolute levels are present at every
  step, so the \citet{arenberg2022} condition is enforced by the arithmetic
  rather than assumed away. With every propensity set to one the extended engine
  reproduces the ordinary one to machine precision, which is what makes any
  difference in the results attributable to the mechanism rather than to a
  second implementation of the bookkeeping.

  \item \textbf{Calibrates it from evidence about something other than the
  answer.} The two mainstream parameters are the dispersion of completed family
  size and the parent--child correlation in it, taken from independent sources
  (\Cref{sec:calibration}). Their combination reproduces an intergenerational
  covariance, which is a weaker and more defensible claim than a heritability,
  though \Cref{sec:response} is explicit that the covariance is constructed from
  two datasets rather than measured in one.

  \item \textbf{Reports a threshold instead of a preferred number.} How the
  fertility \emph{environment} moves --- how many children a person of fixed
  disposition has under prevailing conditions --- cannot be measured in advance,
  and it is where a long-run projection can be made to say anything one likes.
  Because selection and the environment both multiply fertility, we can ask how
  much additional environmental decline would exactly offset measured selection,
  and report that (\Cref{sec:boundary}) rather than choosing a value.
\end{enumerate}

\subsection{What we find}
\label{sec:findings-preview}

Selection is real, measurable, and smaller than its advocates suggest: it shifts
the level of the long-run curve without restoring growth, and it is
overwhelmingly a property of ordinary variation rather than of the visible
high-fertility minorities that dominate discussion of it. Cancelling it takes a
different amount of environmental decline depending on whether the target is the
fertility rate or the accumulated population. \Cref{sec:results} gives the
magnitudes and \Cref{sec:boundary} the thresholds.

\subsection{What this paper does not claim}
\label{sec:not-claims}

\caveatbox{Three things to hold onto while reading}{%
\emph{It is not a forecast.} Every comparison is conditional on a stated path,
and \Cref{sec:boundary} exists because the decisive quantity is unknown. We
report where the answer changes sign, not what the answer is.\\[0.4em]
\emph{It is not a claim about genes.} The transmission parameter is a
parent--child correlation in realised family size. It combines whatever genetic,
material and cultural continuity exists between generations, and nothing in the
argument requires separating them. \citet{hruschka2016} show that much of the
observed variation in completed family size can arise from a counting process
rather than from stable differences between women; \Cref{sec:calibration}
explains what the calibration concedes to that and what survives it.\\[0.4em]
\emph{It is not an economic model.} The fertility environment is a stated
multiplier path, not a projection of income, housing or childcare costs. We
report thresholds in it. We do not forecast it, and we take no position on what
drives it.}
